\section{Formalization}\label{sec:formalization}

The development is carried out in Lean~4 over Mathlib, and a statement
whose head carries the tag \emph{Lean} is verified there.
Tables~\ref{tab:lean} and~\ref{tab:leantwo} name the declaration behind
every tagged statement, and all of them where a statement is assembled from
several.  Three tags are not discharged in full, and \S\ref{sec:boundaries}
records which.

The formal statements are the originals, and the statements in the text
are their transcriptions.

\subsection{What is formalized}\label{sec:whatformal}

The formal definitions are those of \S\ref{sec:first} verbatim.  A design is a
structure carrying the vectors, the weights, positivity of the weights, their
normalization, and \eqref{eq:parseval}; domination of a subset is positive
semidefiniteness of $\SC-I$; and $\Gtz(m,k)$ is the assertion that every
design of size $m$ and rank $k$ admits a dominating subset of cardinality $k$.
No nondegeneracy is assumed: repeated atoms and zero atoms are admitted, and
the results of \S\ref{sec:ties} depend on that.

There is no separate real and complex development.  The structure and the
elementary theory are written once for scalars in an \texttt{RCLike} field
with real weights, and $\R$ and $\C$ are instantiations, definitionally so
over $\R$; Theorems~\ref{thm:corankone} and~\ref{thm:duality} and
Proposition~\ref{prop:sizerank} are proved at that generality and used at both.
Where a statement is genuinely field-dependent, as
Theorem~\ref{thm:complex} and Theorem~\ref{thm:ranktwo} are, it is
proved at the field where it holds.

The chart of
\S\ref{sec:chart} is formalized as the projection $P=VV\adj$ of
\eqref{eq:chart} together with the equivalence of Lemma~\ref{lem:chart}, and
the compactification of \S\ref{sec:compactness} as the predicate of
Definition~\ref{def:chartpoint} on pairs $(P,t)$ with $P$ a symmetric
idempotent of trace $k$ and $t$ in the closed simplex.  The stationarity
system of \S\ref{sec:interior} is a structure bundling
\eqref{eq:statg}--\eqref{eq:statt} with the active family and the
multipliers, and the results of that section are implications from it.

Graphic designs, used for the diamond of Example~\ref{ex:diamond}, are
formalized through \eqref{eq:graphicatom}.  The whitener is then an
existential and the atoms are determined only up to the orthogonal group,
so what is rational, and what the formal statements quantify over, is the
graph data.

\begin{table}[htbp]
\small
\begin{tabular}{@{}ll@{}}
\toprule
Statement & Declaration \\
\midrule
\multicolumn{2}{@{}l}{\emph{\S\ref{sec:first}, elementary cases}}\\
Lemma~\ref{lem:span}           & \texttt{Gtz.rank\_le\_size\_of\_fieldDesign}, \\
                               & \texttt{Gtz.eq\_zero\_of\_forall\_atom\_dot\_eq\_zero} \\
Lemma~\ref{lem:trace}          & \texttt{Gtz.trace\_identity} \\
Lemma~\ref{lem:share}          & \texttt{Gtz.weighted\_leverage\_le\_one} \\
Lemma~\ref{lem:coparseval}     & \texttt{Gtz.coParsevalOperator\_posDef} \\
Proposition~\ref{prop:sizerank}& \texttt{Gtz.gtzWeighted\_square} \\
Proposition~\ref{prop:rankone} & \texttt{Gtz.gtz\_rank\_one} \\
Proposition~\ref{prop:classical}
                               & \texttt{Gtz.gtzOriginal\_of\_projectionCovering}, \\
                               & \texttt{Gtz.gtzOriginal\_iff\_frameProjectionCovering} \\
Lemma~\ref{lem:direction}      & \texttt{Gtz.exists\_atom\_covering\_direction} \\
Example~\ref{ex:tetra}         & \texttt{Gtz.tetraDesign\_isTie} \\
\addlinespace
\multicolumn{2}{@{}l}{\emph{\S\ref{sec:chart}, the projection chart}}\\
Proposition~\ref{prop:chartid} & \texttt{Gtz.projectionChart\_diagonal},
                                 \texttt{Gtz.trace\_projectionChart} \\
Lemma~\ref{lem:congruence}     & \texttt{Gtz.chartBlock\_sub\_weightDiagonal} \\
Lemma~\ref{lem:flip}           & \texttt{Gtz.posSemidef\_one\_sub\_iff\_contraction}, \\
                               & \texttt{Gtz.contraction\_flip} \\
Lemma~\ref{lem:chart}          & \texttt{Gtz.fieldDominates\_iff\_posSemidef\_chartBlock\_finset} \\
Corollary~\ref{cor:minor}      & \texttt{Gtz.det\_projectionBlock\_sub\_weightDiagonal} \\
Lemma~\ref{lem:rayleigh}       & \texttt{Gtz.chartGapRayleigh\_pos\_of\_fixed}, \\
                               & \texttt{Gtz.chartGapRayleigh\_eq\_of\_killed} \\
Remark~\ref{rem:nogo}          & \texttt{Gtz.inertiaNoGoMatrix\_positive\_direction} and companions \\
Lemma~\ref{lem:whiten}         & \texttt{Gtz.exists\_congruence\_to\_one} \\
Theorem~\ref{thm:duality}      & \texttt{Gtz.weighted\_naimark\_duality} \\
Corollary~\ref{cor:dualcell}   & \texttt{Gtz.gtzWeighted\_dual\_iff} \\
Lemma~\ref{lem:rankone}        & \texttt{Gtz.posSemidef\_sub\_vecMulVec\_iff}, \\
                               & \texttt{Gtz.posDef\_sub\_vecMulVec\_iff} \\
Theorem~\ref{thm:corankone}    & \texttt{Gtz.corank\_one\_dominating\_erasure}, \\
                               & \texttt{Gtz.fieldGtzWeighted\_corank\_one},
                                 \texttt{Gtz.isTie\_iff\_forall\_pivot\_eq\_one}, \\
                               & \texttt{Gtz.exists\_realDesign\_exactlyTied},
                                 \texttt{Gtz.exists\_complexDesign\_exactlyTied} \\
\addlinespace
\multicolumn{2}{@{}l}{\emph{\S\ref{sec:complex}, the complex problem}}\\
Proposition~\ref{prop:cpositive}
                               & \texttt{Gtz.complexGtzWeighted\_square}, \\
                               & \texttt{Gtz.complexGtzWeighted\_corank\_one} \\
Example~\ref{ex:sic}           & \texttt{Gtz.sicParseval}, \texttt{Gtz.sicNorm} \\
Lemma~\ref{lem:sicvalue}       & \texttt{Gtz.sicDesign\_valueAtMost},
                                 \texttt{Gtz.pair\_not\_posSemidef} \\
Lemma~\ref{lem:spike}          & \texttt{Gtz.spikePaddedDesign\_valueAtMost} \\
Lemma~\ref{lem:split}          & \texttt{Gtz.splitAtomWeight\_valueAtMost} \\
Theorem~\ref{thm:ladder}       & \texttt{Gtz.complexRankConstantAtMostAtSize\_rank\_add} \\
Theorem~\ref{thm:complex}      & \texttt{Gtz.complexGtzWeighted\_iff\_size\_le\_rank\_add\_one} \\
\quad (negative half)          & \texttt{Gtz.not\_complexGtzWeighted\_of\_rank\_add\_two\_le\_size} \\
Example~\ref{ex:trine}         & \texttt{Gtz.trineBinding\_leastEigenvalue}, \\
                               & \texttt{Gtz.complexRankConstantAtMost\_three\_trine} \\
Example~\ref{ex:hesse}         & \texttt{Gtz.hesseAtom\_leverage},
                                 \texttt{Gtz.hesseAtom\_overlapCube} \\
Proposition~\ref{prop:hesse}   & \texttt{Gtz.hesseMargin\_isRoot},
                                 \texttt{Gtz.hesseDesign\_valueAtMost}, \\
                               & \texttt{Gtz.hesseMargin\_isGreatest} \\
\addlinespace
\multicolumn{2}{@{}l}{\emph{\S\ref{sec:real}, the decided real cells}}\\
Theorem~\ref{thm:light}        & \texttt{Gtz.dominating\_of\_light\_atom} \\
Corollary~\ref{cor:heavy}      & \texttt{Gtz.gtzWeightedAll\_of\_heavy} \\
Theorem~\ref{thm:ranktwo}      & \texttt{Gtz.gtz\_rank\_two} \\
Theorem~\ref{thm:coranktwo}    & \texttt{Gtz.gtzWeighted\_corank\_two} \\
Corollary~\ref{cor:corank}     & \texttt{Gtz.gtzWeighted\_of\_le\_five} \\
Lemma~\ref{lem:null}           & \texttt{Gtz.exists\_null\_direction} \\
Theorem~\ref{thm:reduce}       & \texttt{Gtz.exists\_reduced\_design} \\
Theorem~\ref{thm:window}       & \texttt{Gtz.crystallization} \\
Lemma~\ref{lem:repeat}         & \texttt{Gtz.not\_dominates\_of\_repeated\_atom} \\
Proposition~\ref{prop:splitting}& \texttt{Gtz.gtzWeighted\_of\_succ} \\
Theorem~\ref{thm:topcell}      & \texttt{Gtz.gtzWeightedAll\_of\_top}, \\
                               & \texttt{Gtz.gtz\_iff\_top\_of\_each\_rank} \\
Proposition~\ref{prop:spike}   & \texttt{Gtz.gtzWeighted\_of\_spike} \\
\bottomrule
\end{tabular}
\medskip
\caption{Correspondence between the numbered statements and the formal
development, first part: the elementary theory, the chart, the complex answer
and the decided real cells.}
\label{tab:lean}
\end{table}

\begin{table}[htbp]
\small
\begin{tabular}{@{}ll@{}}
\toprule
Statement & Declaration \\
\midrule
\multicolumn{2}{@{}l}{\emph{\S\ref{sec:ties}, the equality locus}}\\
Theorem~\ref{thm:tieseverywhere}& \texttt{Gtz.exists\_isTie\_of\_weights\_of\_classes}, \\
                               & \texttt{Gtz.splitClassDesign\_isTie},
                                 \texttt{Gtz.splitTetraDesign\_isTie}, \\
                               & \texttt{Gtz.splitSevenDesign\_isTie} \\
Example~\ref{ex:split}         & \texttt{Gtz.splitTetraDesign\_dominates}, \\
                               & \texttt{Gtz.splitTetraDesign\_no\_strictDominator} \\
Corollary~\ref{cor:nomargin}   & \texttt{Gtz.not\_exists\_card\_posSemidef\_sub\_smul\_one\_of\_isTie}, \\
                               & \texttt{Gtz.not\_gtzWeightedFloor\_rankThree\_of\_one\_lt} \\
Example~\ref{ex:diamond}       & \texttt{Gtz.diamondDesign\_isTie}$^{a}$, \\
                               & \texttt{Gtz.not\_hasParallelPair\_diamondDesign} \\
Corollary~\ref{cor:noparallel} & \texttt{Gtz.not\_hingeHoldsAtSize\_five\_three},
                                 \texttt{Gtz.diamondEdgeVector\_eq} \\
Theorem~\ref{thm:leverageaverage}
                               & \texttt{Gtz.posSemidef\_leverageWeightedAtomSum\_sub\_one} \\
Proposition~\ref{prop:jensen}  & \texttt{Gtz.sq\_rank\_le\_expectedElementary\_one} \\
Proposition~\ref{prop:qonezero}& \texttt{Gtz.mixedCharPoly\_eval\_one\_eq\_zero\_of\_totalTieSupported} \\
\addlinespace
\multicolumn{2}{@{}l}{\emph{\S\ref{sec:compactness}, the variational principle}}\\
Theorem~\ref{thm:boundary}     & \texttt{Gtz.exists\_chartDominates\_of\_weight\_eq\_zero} \\
Corollary~\ref{cor:interiorcex}& \texttt{Gtz.weight\_pos\_of\_forall\_not\_chartDominates} \\
Theorem~\ref{thm:sharp}        & \texttt{Gtz.not\_eventualCompressionLiftClaim} \\
\addlinespace
\multicolumn{2}{@{}l}{\emph{\S\ref{sec:interior}, interior critical points}}\\
Proposition~\ref{prop:quadric} & \texttt{Gtz.quadricLaw\_of\_isQuadricStationaryData}, \\
                               & \texttt{Gtz.coverageLaw\_of\_isQuadricStationaryData}, \\
                               & \texttt{Gtz.multiplierMatrix\_eq\_of\_isQuadricStationaryData} \\
Lemma~\ref{lem:overlap}        & \texttt{Gtz.tightOverlap\_sum\_eq\_one\_of\_isQuadricStationaryData}, \\
                               & \texttt{Gtz.one\_le\_leverage\_of\_isQuadricStationaryData}, \\
                               & \texttt{Gtz.exists\_rayleighProbe\_ge\_rank\_of\_isQuadricStationaryData}, \\
                               & \texttt{Gtz.exists\_mem\_activeSubset\_of\_isQuadricStationaryData} \\
Corollary~\ref{cor:singleactive}
                               & \texttt{Gtz.value\_eq\_rank\_of\_singleActive}, \\
                               & \texttt{Gtz.not\_isQuadricStationaryData\_of\_singleActive\_of\_value\_lt\_one} \\
Proposition~\ref{prop:saturated}& \texttt{Gtz.weight\_mul\_value\_eq\_one\_of\_saturatedAtom}, \\
                               & \texttt{Gtz.one\_le\_value\_of\_saturatedAtom} \\
\addlinespace
\multicolumn{2}{@{}l}{\emph{Appendix~\ref{app:graveyard}, excluded strategies}}\\
Proposition~\ref{prop:noaggregate}
                               & \texttt{Gtz.not\_hasUniformTieAggregateSeven}, \\
                               & \texttt{Gtz.splitSevenDesign\_weightedTieAggregate\_nonpos}, \\
                               & \texttt{Gtz.closedTieFailure\_inhabited\_seven} \\
Proposition~\ref{prop:twomoment}& \texttt{Gtz.dominates\_of\_twoMomentGap\_nonneg} \\
Proposition~\ref{prop:icosa}   & \texttt{Gtz.icosa\_twoMomentGap\_eq},
                                 \texttt{Gtz.icosa\_no\_twoMoment\_certificate}, \\
                               & \texttt{Gtz.icosaDesign\_strictly\_dominates} \\
\bottomrule
\end{tabular}
\medskip
\caption{Correspondence, second part: the equality locus, the variational
principle, the interior system and the excluded strategies.  The declaration
marked $a$ is proved but does not appear in the axiom ledger described in
\S\ref{sec:axioms}.}
\label{tab:leantwo}
\end{table}

\subsection{Axiom discipline}\label{sec:axioms}

A ledger module lists, for the declarations of the development, the axioms on
which each depends, printed by the kernel rather than asserted by the author.
The permitted set is Mathlib's
$\{\texttt{propext},\ \texttt{Classical.choice},\ \texttt{Quot.sound}\}$, and
nothing outside it occurs.  Classical logic is used freely, and nothing in
the development is a constructive claim.

In particular no declaration depends on \texttt{sorryAx}: the development
contains no incomplete proof.  It contains no \texttt{axiom} declaration.
Closed arithmetic over concrete finite index sets is discharged by kernel
evaluation of decidable propositions, which introduces no axiom.  No proof
is closed by compiled evaluation.

The ledger covers exactly what it lists: one
declaration, the entry marked $a$ in Table~\ref{tab:leantwo}, is absent
from it, and is proved in a module duplicating
one that the ledger does cover.  Its proof is checked when that module is
elaborated, but its axiom use is not certified by the ledger.

No numerical assertion in the paper rests on floating-point evaluation.  The
witnesses of \S\ref{sec:sharp} and Appendix~\ref{app:graveyard} are exact
rationals in the formal statements and are checked as such; the icosahedral and
graphic designs of Appendix~\ref{app:data} carry radicals, and the
formal statements use the algebraic relations \eqref{eq:icosaconst} and
\eqref{eq:graphicleverage} that eliminate them.  Remark~\ref{rem:hygiene} records the
conditioning phenomenon that makes this discipline necessary rather than
decorative.

\subsection{Where the formalization stops}\label{sec:boundaries}

Two steps of the argument are not formal.  Both lie in
\S\ref{sec:compactness}--\S\ref{sec:interior}.

The first is the attainment in Theorem~\ref{thm:variational}: a continuous
function on a compact space attains its minimum.  The space of chart points is
compact --- a product of a Grassmannian and a closed simplex --- and the
objective $G$ of Definition~\ref{def:chartpoint} is continuous, being a maximum
of finitely many least-eigenvalue functions of continuous matrix-valued maps.
We use this standard fact without formalizing it.  Its two nontrivial
inputs, Theorem~\ref{thm:boundary} and Corollary~\ref{cor:interiorcex},
which place the minimizer at a design rather than at a boundary point,
are formal.

The second is the passage from a local minimum to the system
\eqref{eq:statg}--\eqref{eq:statt}.  Its ingredients are the spectraplex,
the composition rule and the maximum rule, assembled in \S\ref{sec:kkt};
each is classical
\cite{Danskin1967,Overton1988,HiriartUrrutyYe1995,Clarke1990} and none is
in the ambient library.  So \S\ref{sec:interior} takes the system as its
hypothesis and formalizes its consequences
(Proposition~\ref{prop:quadric}, Proposition~\ref{prop:e1} and the
exclusions) rather than deriving the system from minimality.

The two conjectural inputs of Theorem~\ref{thm:rankthree} speak about the
system, not about minimality, so a proof of either closes the corresponding
cell only together with the two steps above.

Three tags are narrower than the statement they head, and in each case the
gap is one step of the printed proof.
\begin{enumerate}
\item[\rm(i)] Proposition~\ref{prop:rankfloor} is stated over $\FF$ and
  formalized over $\R$, as\\
  \texttt{Gtz.gtzWeightedFloor\_inv\_rank}.  Its proof does not consult the
  field, so the complex case is a transcription and not a mathematical gap.
\item[\rm(ii)] In Proposition~\ref{prop:tielawcorank} the per-atom
  equivalence is formal, as\\
  \texttt{Gtz.jensenRatio\_eq\_inv\_rank\_iff\_leverageIdentity},\\
  and the summation that concludes $m=k+1$ is not.
\item[\rm(iii)] In Lemma~\ref{lem:defect} the two consequences drawn from the
  identity are formal, as\\
  \texttt{Gtz.weight\_mul\_value\_le\_one\_of\_isQuadricStationaryData} and\\
  \texttt{Gtz.weight\_mul\_value\_eq\_one\_of\_saturatedAtom};\\
  the identity itself, one subtraction of \eqref{eq:coverage} from
  \eqref{eq:overlapone}, is not stated on its own.
\end{enumerate}
